% !PRINTED 349
\weberpage{349}

Die Gleichung (8) giebt nun selbst schon einen ziemlich guten Näherungswerth für die reelle Wurzel

\[
x = \sqrt[8]{82} = 1{,}73471.
\]

Ein genauerer Werth, den man erhält, wenn man noch zwei Schritte weiter geht, ist 1,73469, der in der fünften Decimale noch richtig ist.

Eine Einschliessung der gesuchten Wurzel in zwei Grenzen giebt diese Methode nicht. Das Kennzeichen, ob ein genügender Grad von Genauigkeit erreicht ist, besteht darin, dass, wenn $s_\mu$, $s_{\mu'}$ zwei auf einander folgende, zur Berechnung benutzte Potenzsummen der Wurzeln sind, diese sich mit genügender Genauigkeit verhalten, wie die $\mu^{\text{te}}$ zur $\mu'^{\text{ten}}$ Potenz einer Grösse, oder dass annähernd

\[
\mu \log s_{\mu'} - \mu' \log s_\mu = 0
\]

ist. Man muss aber diese Prüfung bei mehreren auf einander folgenden Gliedern vornehmen, da in besonderen Fällen diese Relation scheinbar erfüllt sein kann und bei der späteren Rechnung wieder aufhört.

\begin{center}
\S.\ 112.

\gesperrt{Trigonometrische Auflösung cubischer}

\gesperrt{Gleichungen.}
\end{center}

Wir besprechen nun noch einige auf Gleichungen von speciellen Formen anwendbare Methoden der numerischen Auflösung. Das Ziel dieser Methoden ist, die allgemein verbreiteten Tafeln der trigonometrischen Functionen und der Logarithmen für die Auflösung von Gleichungen nutzbar zu machen. Wir wenden uns zunächst zur Betrachtung der cubischen Gleichungen, die wir immer in der reducirten Form

\weq{1}{x^3 + a x + b = 0}

annehmen, worin $a$ und $b$ reelle Zahlen sind. Da die Vertauschung von $x$ mit $-x$ gleichbedeutend ist mit der Vertauschung von $b$ mit $-b$, so können wir uns auf die Annahme beschränken, dass $b$ negativ sei, und es bleiben dann noch drei verschiedene Fälle zu betrachten. Wir setzen $b = -g$ und, je nachdem $a$ positiv oder negativ ist, $a = \pm e$. Dann haben wir
%CONT>>
